\documentclass[11pt]{article}
\usepackage[T1]{fontenc}
\usepackage[utf8]{inputenc}
\usepackage[letterpaper,margin=1in]{geometry}
\usepackage{amsmath,amsthm,mathtools}
\usepackage{newtxtext,newtxmath}
\usepackage{microtype,graphicx,aliascnt,flafter}
\usepackage{tikz}
\usetikzlibrary{arrows.meta,positioning,fit,calc}
\usepackage[hidelinks]{hyperref}
\usepackage[nameinlink,noabbrev]{cleveref}
\usepackage{fancyhdr}
\setlength{\headheight}{14pt}
\pagestyle{fancy}\fancyhf{}
\fancyhead[L]{An I/O lower bound for exact attention}
\fancyfoot[C]{\thepage}
\renewcommand{\headrulewidth}{0.3pt}
\setlength{\parindent}{1.25em}
\setlength{\parskip}{0pt}
\setlength{\emergencystretch}{1.5em}
\allowdisplaybreaks[1]
\widowpenalty=10000
\clubpenalty=10000
\newtheorem{theorem}{Theorem}[section]
\newaliascnt{lemma}{theorem}
\newtheorem{lemma}[lemma]{Lemma}
\aliascntresetthe{lemma}
\newaliascnt{proposition}{theorem}
\newtheorem{proposition}[proposition]{Proposition}
\aliascntresetthe{proposition}
\newaliascnt{corollary}{theorem}
\newtheorem{corollary}[corollary]{Corollary}
\aliascntresetthe{corollary}
\theoremstyle{definition}
\newaliascnt{definition}{theorem}
\newtheorem{definition}[definition]{Definition}
\aliascntresetthe{definition}
\theoremstyle{remark}
\newaliascnt{remark}{theorem}
\newtheorem{remark}[remark]{Remark}
\aliascntresetthe{remark}
\crefname{lemma}{Lemma}{Lemmas}
\crefname{theorem}{Theorem}{Theorems}
\crefname{proposition}{Proposition}{Propositions}
\crefname{corollary}{Corollary}{Corollaries}
\crefname{definition}{Definition}{Definitions}
\crefname{remark}{Remark}{Remarks}
\crefname{section}{Section}{Sections}
\crefname{appendix}{Appendix}{Appendices}
\newcommand{\R}{\mathbb R}
\newcommand{\F}{\mathbb F}
\newcommand{\cH}{\mathcal H}
\newcommand{\cT}{\mathcal T}
\newcommand{\cE}{\mathcal E}
\newcommand{\trdeg}{\operatorname{trdeg}}
\newcommand{\rank}{\operatorname{rank}}
\newcommand{\softmax}{\operatorname{softmax}}
\newcommand{\diag}{\operatorname{diag}}
\newcommand{\IO}{\operatorname{IO}}
\newcommand{\vecop}{\operatorname{vec}}
\newcommand{\dd}{\mathrm d}
\hypersetup{pdftitle={An I/O Lower Bound for Exact Attention},pdfauthor={Samuel Mausberg},pdfsubject={Exact attention in an analytic real-arithmetic memory model}}
\title{An I/O Lower Bound for Exact Attention}
\author{Samuel Mausberg\\Independent Researcher}
\date{}

% Keep all prose, including captions and references, at the 11-point body size.
\crefname{equation}{equation}{equations}
\Crefname{equation}{Equation}{Equations}
\crefname{figure}{Figure}{Figures}
\Crefname{figure}{Figure}{Figures}
\setlength{\textfloatsep}{16pt plus 2pt minus 2pt}
\setlength{\floatsep}{12pt plus 2pt minus 2pt}
\setlength{\intextsep}{14pt plus 2pt minus 2pt}
\begin{document}
\hypersetup{pageanchor=false}
\begin{titlepage}
\thispagestyle{empty}
\vspace*{0.3in}
\begin{center}
{\LARGE An I/O Lower Bound for Exact Attention\par}
\vspace{0.25in}
{\large Samuel Mausberg\par}
\vspace{0.05in}
Independent Researcher
\end{center}
\vspace{0.35in}
\begin{center}Abstract\end{center}
Every bounded deterministic real-arithmetic program that computes exact
softmax attention on $Q,K,V\in\R^{n\times d}$ requires
$\Omega(nd+n^2/M)$ transfers between slow memory and an $M$-word cache,
for $n\ge d^2$, $d\ge2$, and $M\ge d^2$. The allowed operations are
$+,-,\times,\div,\exp$; recomputation and finite branching are permitted.
No pairwise intermediate value is prescribed. The proof adjoins the
entries of each epoch's local Jacobian to a field that retains all earlier
such entries. An epoch contributes at most $4M^2$ generators, whereas
the output derivatives determine $n(n-1)$ algebraically independent
kernel ratios. The bound is tight for each fixed $d$. We also show why
field containment does not recover the uniform dependence on $d$.
For powers of two $d\ge8$ dividing $n$, a Strassen-based prefix reaches
the entire attention coefficient field after
$O(nd^{\sigma-1}+n^2d^{\sigma-2}/M)$ transfers, where $\sigma=\log_2 7$.
This prefix can occur in a correct $O(n^2d)$-work program, yet on a
parameter sequence it costs $o(nd+n^2d/M)$. It computes an auxiliary
vector, not attention. A separate numerical compression theorem gives
$\Omega(nd+n^2d^2/M)$ when the epochs recover all pair scores or kernels;
that recovery hypothesis is not part of the unrestricted lower bound.
\end{titlepage}
\hypersetup{pageanchor=true}
\setcounter{page}{1}

\section{Introduction}
For $Q,K,V\in\R^{n\times d}$, exact attention is the map
\begin{equation}
 W_{ij}=\exp(q_i\cdot k_j),\qquad
 A_{ij}=\frac{W_{ij}}{\sum_{h=1}^n W_{ih}},\qquad Y=AV,
 \label{eq:attention}
\end{equation}
where $q_i$ and $k_j$ are rows. The inputs begin in slow memory, the
output must be written there, and arithmetic operates on at most $M$
fast real words. We count scalar loads and stores.

The blocking schedule underlying FlashAttention uses
$O(nd+n^2d^2/M)$ transfers and $O(n^2d)$ arithmetic work
\cite{dao2022}. A formula with $n^2$ pairwise scores does not require
a program to form those scores. Products of coordinate-wise
exponentials already avoid them, and more general rearrangements
can combine terms before an individual kernel entry is available.
Our lower bound makes no assumption about these intermediate values.

\begin{theorem}[Exact attention]
\label{thm:attention}
For $n\ge d^2$, $d\ge2$, and $M\ge d^2$, every bounded deterministic
program over $+,-,\times,\div,\exp$ in the model of
\cref{sec:model} that computes \cref{eq:attention} on all real inputs
has worst-case transfer complexity
\begin{equation}
 \IO\ge\frac1{32}\left(nd+\frac{n^2}{M}\right).
 \label{eq:main-lower}
\end{equation}
The program may recompute values, branch on sign comparisons, and use
fast matrix multiplication. It needs only a finite worst-case
instruction bound for each fixed $(n,d,M)$.
\end{theorem}

For each fixed $d$, this matches the blocking upper bound, with
comparison constants depending on that dimension. Uniformly in $d$,
the bounds remain
\begin{equation}
 \Omega\!\left(nd+\frac{n^2}{M}\right)
 \ \le\ \IO^*(n,d,M)\ \le\
 O\!\left(nd+\frac{n^2d^2}{M}\right),
 \label{eq:bracket}
\end{equation}
where $\IO^*$ is the least worst-case cost among exact families with
$O(n^2d)$ work. The lower bound applies more broadly, as the theorem
states. We do not prove even the intermediate
$\Omega(nd+n^2d/M)$ bound when $d$ grows. The term $n^2/M$ exceeds the
compulsory $nd$ only when $n>Md$, so the lower bound is informative in
the long-context regime.

The proof follows derivatives through the memory schedule. Since
$Y$ is linear in $V$, its Jacobian contains every entry of $A$.
The ratios $A_{ij}/A_{in}$ are $n(n-1)$ algebraically independent
functions over the rational input field. An epoch of at most $M$
transfers has at most $2M$ incoming and $2M$ outgoing values, so its
local Jacobian has at most $4M^2$ entries. The chain rule puts all
output derivatives in the field generated by these local entries
throughout the execution. This requires $\Omega(n^2/M^2)$ epochs
when the nontrivial term dominates. Derivatives are proof objects;
the program never has to compute them.

Using field containment as the sole completion condition does not
recover the dimension factor. Write
\begin{equation}
 F_0=\R(Q,K,V),\qquad
 R_{ij}=\exp\bigl(q_i\cdot(k_j-k_n)\bigr)\quad(j<n),\qquad
 \cT=F_0(R_{ij}:i\in[n],\ j<n).
 \label{eq:target}
\end{equation}
For powers of two $d\ge8$ dividing $n$, we construct a prefix whose
local-Jacobian history already contains $\cT$, at a cost of
\begin{equation}
 O\!\left(nd^{\sigma-1}+\frac{n^2d^{\sigma-2}}{M}\right),
 \qquad \sigma=\log_2 7.
 \label{eq:prefix-intro}
\end{equation}
It uses Strassen's bilinear decomposition to compute weighted sums of
scalar exponentials. Differentiation separates the summands in the
function field, where integer products recover the original kernels.
Those products are not evaluated by the prefix. At $n=d^3$ and
$M=d^2$, its cost is $\Theta(d^{\sigma+2})=o(d^5)$, while
$nd+n^2d/M=\Theta(d^5)$. Appending a correct attention computation
gives a uniform $O(n^2d)$-work program with the same cheap prefix.
Thus containing the coefficient field cannot, on its own, certify
that the larger transfer cost has already been paid.

Numerical recovery imposes a different requirement. If $m$
differentiable real summaries can decode $B$ selected pair scores or
kernels, then $B\le m^2/d^2+2m$. \Cref{app:recovery} proves this
compression bound and the full attention lower bound under pair
coverage. It is a conditional statement about numerical values,
not a consequence of field containment.

\paragraph{Prior work.}
Hong and Kung's red-blue pebble game relates I/O to small-boundary
partitions of a specified computation graph
\cite[Theorem~3.1]{hongkung1981}. We partition the executed program,
without prescribing its graph. Saha and Ye permit fast matrix
multiplication but explicitly assume that the program computes
$QK^{\mathsf T}$, though it need not store that matrix
\cite[Section~4]{sahaye2024}. With $\mathbb F_q$ words and $q>n$, their
Theorem~4.8 gives $\Omega(\min\{n^2d^2/M,n^2\})$ transfers. These
numbers refer to arXiv:2402.07443v1. The present model instead allows
arbitrary real words and analytic arithmetic, but no digit extraction
or real-to-integer conversion. Their bound is stronger in $d$ but
applies to programs that compute $QK^{\mathsf T}$. To our knowledge,
\cref{thm:attention} is the first I/O lower bound for exact attention
beyond $\Omega(nd)$ that does not prescribe the computation of scores or
kernels.
The blocking comparison is the upper
bound of \cite[Theorem~2]{dao2022}; \cref{app:upper} supplies a schedule
for the full parameter range in \cref{eq:bracket}.

Abelson used derivatives to lower-bound communication between
processors \cite{abelson1980}. The numerical factorisation argument
in \cref{lem:factor-rank} uses this chain-rule principle. Grigoriev
\cite{grigoriev2008} and Blanc, Dwivedi, Hansen, Limaye, and Mahajan
\cite{blanc2026} prove lower bounds for two-party protocols that
exchange values of polynomials in real inputs. Here we count generators
contributed by local Jacobians throughout an execution, and the input
is not split between parties. The extended Baur--Strassen theorem of Saha, Xu, Ye, and Yu
bounds differentiation in arithmetic circuits with exponentials and
logarithms \cite[Theorem~4.4]{sahaxuyeyu2026}; that circuit-size
statement does not by itself preserve an I/O schedule. Our argument
uses no numerical differentiation algorithm.

\paragraph{Formalization.}
A Lean~4 formalization accompanies this paper. It proves
\cref{lem:independence,lem:open,lem:factor-rank,lem:pair-rank,thm:compression},
the identities of \cref{lem:strassen}, and the independence and derivative
identities behind \cref{prop:attention-field,thm:bilinear-prefix}. It also
proves \cref{thm:history} for execution traces and derives
\cref{eq:main-lower} for every trace whose outputs carry the Jacobian of
attention. The passage from a program to its trace and the transfer counts
of the constructions are not formalized.

\section{The machine and its open execution paths}
\label{sec:model}
For each integer triple $(n,d,M)$, a machine has $M$ fast real registers
and an unbounded array of slow real words. The input matrices occupy
designated slow addresses. Output addresses are separate and initially
unwritten. Initial register values and program constants are
independent of the inputs. Loading or storing one scalar costs one
transfer. Arithmetic acts only on fast registers and may overwrite an
operand. 

The scalar instructions are $+,-,\times,\div,\exp$, copies, and sign
comparisons of register values or their differences. A division must
have a nonzero denominator on its executed branch. Addresses use
discrete counters and finite control. There is no instruction that
converts an arbitrary real value to an integer, extracts a digit,
performs a real-indexed lookup, or evaluates a function purportedly
encoded in a real word. Code constants do not provide such operations.
Recomputation and fast matrix multiplication are permitted.

Programs are deterministic and have a finite worst-case instruction
bound at each fixed parameter triple. All instructions, including
comparisons and counter updates, count as work. A finite I/O claim also
has a finite worst-case transfer bound. Thus the complete computation
can be unrolled into a finite decision tree. The lower bounds allow
any finite work bound and arbitrary input-independent real constants.
Algorithmic upper bounds use $O(n^2d)$ work with a constant uniform in
the parameters. All complexity statements are worst-case statements
on arbitrary real inputs. The same conventions apply to a general
finite input vector $x$ and output vector $f(x)$.

\begin{lemma}[An open execution path]
\label{lem:open}
A bounded program defined on a nonempty open input set has one
execution path followed on a nonempty open subset. On that subset,
all values appearing in the path are analytic functions of the inputs,
and the accessed addresses are fixed.
\end{lemma}
\begin{proof}
Start with a connected open subset. Arithmetic preserves analyticity
where the denominators of executed divisions are nonzero. At a sign
test, the analytic test function either vanishes identically or is
nonzero on some smaller open subset, on which one sign is fixed.
Choose the corresponding branch. Counter values and addresses are
then fixed along the path. Repeating this step reaches a leaf because
the decision tree is finite. The resulting nonempty open set may be
shrunk to be connected.
\end{proof}

Fix such a connected open set $U$. Analytic functions on $U$ form an
integral domain; we work in its fraction field. The rational function
field $\R(x_1,\ldots,x_p)$ embeds in this field because no nonzero
polynomial vanishes on an open set. Statements about algebraic
independence below concern functions on $U$, rather than particular
real numbers. Poles introduced by a field identity have no implication
for the instructions executed by the program.

For an analytic vector $f(x)$, define
\begin{equation}
 \tau_1(f)=\trdeg_{\R(x)}\R(x)(Df).
 \label{eq:tau}
\end{equation}
The right side is the transcendence degree of the field generated by
the scalar entries of the Jacobian. It can exceed the number of input
coordinates. Its numerical rank cannot exceed the number of inputs.

\begin{lemma}[Numerical factorisation]
\label{lem:factor-rank}
Suppose $C:U\to\R^m$ and $D$ are continuously differentiable, with $D$
defined on a neighbourhood of $C(U)$. If $F=D\circ C$, then
$\rank DF(x)\le m$ on $U$. In particular, fewer than $p$ such numerical
summaries cannot encode and decode all $p$ independent real input
coordinates on an open set.
\end{lemma}
\begin{proof}
The chain rule gives $DF(x)=DD(C(x))\,DC(x)$. Its rank is at most $m$.
Apply this to $F(x)=x$ for the last statement.
\end{proof}

On an open execution path, this rules out lossless packing with
numerical differentiable decoding. Field membership supplies no such
decoder: it is an identity between functions, not an instruction for
recovering their numerical values.

\section{The derivative field of an execution}
\label{sec:history}
An epoch is a consecutive part of an execution with at most $M$
transfers. Partition a path with $I$ transfers into $e$ epochs, using
$M$ transfers per epoch except possibly the last. Arithmetic before the
first transfer and after the last is included. A path with no transfers
has one epoch. In all cases,
\begin{equation}
 e\le I/M+1.
 \label{eq:epochs}
\end{equation}

\begin{lemma}[Boundary factorisation]
\label{lem:boundary}
For each epoch $t$, there are incoming values $C_t(x)\in\R^{a_t}$ and
outgoing values $B_t(x)\in\R^{b_t}$, with $a_t,b_t\le2M$, and an
analytic map $\Phi_t$ on a neighbourhood of a realised boundary point
such that
\begin{equation}
 B_t(x)=\Phi_t(C_t(x)).
 \label{eq:boundary}
\end{equation}
Every newly stored slow word and every final cache value occurs in
$B_t$. Incoming values are initial cache values or versions of slow
words that existed before the epoch.
\end{lemma}
\begin{proof}
Trace the epoch with its addresses fixed. Its initial cache contributes
at most $M$ arguments. For a load from a word whose current version
predates the epoch, add its value as an external argument. There are
at most $M$ such loads. A load that follows a same-epoch store is
replaced by the expression for that stored version, including through
any intervening overwrites. It introduces no external argument.

Every instruction is now an expression in these arguments. Record
all store values, at most $M$, and the final cache, at most $M$, as the
outgoing vector. At any realised boundary point the original executed
denominators are nonzero. They remain nonzero in a sufficiently small
neighbourhood of that point when the incoming arguments are treated
as independent. The resulting expressions define $\Phi_t$ there.
Shrinking the input open set for the finitely many epochs gives the
claimed local factorisations. They need not be unique off the realised
set of boundary values.
\end{proof}

\begin{figure}[t]
 \centering\begin{tikzpicture}[font=\normalsize,>=Latex,
 box/.style={draw,align=center,inner sep=7pt,minimum height=1.9cm}]
 \node[box,text width=3.15cm] (in)
 {$C_t(x)$\\[4pt]initial cache and\\external loads\\$a_t\le2M$};
 \node[box,right=0.7cm of in,text width=4.1cm] (local)
 {$B_t=\Phi_t(C_t)$\\[4pt]at most $M$ transfers\\same-epoch reloads\\substituted internally};
 \node[box,right=0.7cm of local,text width=3.15cm] (out)
 {$B_t(x)$\\[4pt]all stores and\\final cache\\$b_t\le2M$};
 \draw[->] (in.east)--(local.west);
 \draw[->] (local.east)--(out.west);
 \node[below=0.65cm of local,align=center] (jac)
 {$\Lambda_t=D\Phi_t(C_t)$: at most $a_tb_t\le4M^2$ entries\\[6pt]
 $\mathcal H_t=\mathcal H_{t-1}(\Lambda_t)$\qquad $DB_t=\Lambda_t\,DC_t$};
 \draw[->,dashed] (local.south)--(jac.north);
\end{tikzpicture}

 \caption{One epoch and its local Jacobian. The history retains the
 entries of $\Lambda_t$ as functions, without making their numerical
 values available to the machine. A reloaded word inherits its derivatives from earlier epochs.}
 \label{fig:epoch}
\end{figure}

As illustrated in \cref{fig:epoch}, use the maps supplied by the
executed expressions in \cref{lem:boundary} and set
\begin{equation}
 \Lambda_t(x)=D\Phi_t(C_t(x)),\qquad
 \cH_0=F_0=\R(x),\qquad
 \cH_t=\cH_{t-1}(\Lambda_t).
 \label{eq:history}
\end{equation}
Parentheses mean adjoining all scalar entries. These are ordinary
fields. No further differentiation is implicit in their definition.

\begin{theorem}[I/O and the first-derivative field]
\label{thm:history}
Let a bounded analytic program compute $f(x)$ on a nonempty open set,
using $M$ fast words and $I$ transfers on an open execution path. Then
\begin{equation}
 \tau_1(f)\le\sum_{t=1}^e a_tb_t\le4M^2e,
 \qquad
 I\ge\frac{\tau_1(f)}{4M}-M.
 \label{eq:history-bound}
\end{equation}
\end{theorem}
\begin{proof}
The induction tracks every extant word in both memories. A later
load therefore reuses the derivatives of a stored word rather than
introducing a new input function.

Initially, first derivatives of input words and constants are $0$ or
$1$, hence belong to $F_0$. Suppose that every first derivative of
every extant word after epoch $t-1$ belongs to $\cH_{t-1}$. By
\cref{lem:boundary}, this holds for all entries of $DC_t$. Differentiating
\cref{eq:boundary} gives
\begin{equation}
 DB_t=\Lambda_t\,DC_t.
 \label{eq:chain}
\end{equation}
Every scalar entry on the right belongs to $\cH_t$. Newly stored words
and final cache values occur in $B_t$; an unchanged slow word keeps
its previous derivatives. Overwritten versions need not be kept
as numerical words, although their previously adjoined functions
remain in the history field. The induction follows.

All output derivatives therefore belong to $\cH_e$. Adjoining $a_tb_t$
elements increases transcendence degree by at most $a_tb_t$, regardless
of any algebraic dependencies among them. Consequently
\[
 \tau_1(f)\le\trdeg_{F_0}\cH_e
 \le\sum_{t=1}^e a_tb_t\le4M^2e.
\]
Substitute \cref{eq:epochs} and rearrange.
\end{proof}

A reload with no further arithmetic has an identity local map and contributes
no new generator. A long computation inside an epoch may produce many
exponentials, but only the entries of the epoch's boundary Jacobian
are adjoined in \cref{eq:history}. Closure under further differentiation would give a different and
invalid charge. For example, $H_L(a,x)=\sum_{\ell=1}^L e^{a^\ell x}$
uses two loads, one store, and a constant-size cache. Its successive
$x$-derivatives recover all $L$ exponentials by a Vandermonde system,
but its first-derivative field has only two generators. Thus an epoch
may not be charged for every derivative of a cached function.

\section{The attention coefficient field}
\label{sec:attention}
\begin{lemma}[Exponentials of polynomials]
\label{lem:independence}
Let $p_1,\ldots,p_s\in\R[x]$. Suppose every nonzero integer linear
combination of them is nonconstant. Then $e^{p_1},\ldots,e^{p_s}$ are
algebraically independent over $\R(x)$ on every nonempty open set.
\end{lemma}
\begin{proof}
Clear denominators in a putative polynomial relation among the
exponentials and collect monomials to obtain
\begin{equation}
 \sum_{a\in S} r_a(x)e^{P_a(x)}=0,
 \qquad P_a=\sum_{j=1}^s a_jp_j,
 \label{eq:exponential-relation}
\end{equation}
where $S$ is a finite set of distinct nonnegative integer vectors and
each $r_a$ is a nonzero polynomial. The analytic identity extends from
an open set to all of $\R^p$. For a generic vector $z$, restriction to
$x=uz$ keeps every $r_a(uz)$ nonzero as a polynomial in $u$ and makes
every difference $P_a(uz)-P_b(uz)$ nonconstant. Such $z$ exists: each
excluded condition forces a nonzero homogeneous component of one of
finitely many polynomials to vanish.

For sufficiently large positive $u$, the distinct exponent polynomials
are strictly ordered. Divide \cref{eq:exponential-relation} by the
largest exponential. Each remaining term is a polynomial times
$e^{-q(u)}$, where $q(u)\to+\infty$ is a nonconstant polynomial. These
terms tend to zero faster than every inverse power of $u$. The
remaining coefficient is a nonzero polynomial and cannot do so.
\end{proof}

\begin{proposition}[The coefficient field]
\label{prop:attention-field}
For $n\ge2$ and $d\ge1$, the functions $R_{ij}$ of
\cref{eq:target} are algebraically independent over $F_0$, and
\begin{equation}
 F_0(DY)=\cT,\qquad \tau_1(Y)=n(n-1).
 \label{eq:attention-field}
\end{equation}
\end{proposition}
\begin{proof}
The polynomial $q_i\cdot(k_j-k_n)$ has coefficient one on
$Q_{i1}K_{j1}$ for $j<n$. No other polynomial in this family contains
that monomial. Their nonzero integer combinations are therefore
nonconstant, and \cref{lem:independence} applies.

Linearity in $V$ gives $\partial Y_{i1}/\partial V_{j1}=A_{ij}$.
Thus $R_{ij}=A_{ij}/A_{in}$ belongs to $F_0(DY)$, proving one
containment. Conversely, set $R_{in}=1$. The entries of $Y$ are rational
functions of $R$ and $V$, with denominator $1+\sum_{j<n}R_{ij}$.
Every input derivative of $R_{ij}$ is a polynomial times $R_{ij}$.
The field $\cT$ is consequently closed under input differentiation,
and all entries of $DY$ belong to it. Independence gives its
transcendence degree.
\end{proof}

\begin{proof}[Proof of \cref{thm:attention}]
Choose an open execution path using \cref{lem:open}. The function
computed there is still attention, so \cref{thm:history,prop:attention-field}
give
\[
 I\ge\frac{n(n-1)}{4M}-M.
\]
There are also $nd$ compulsory output stores on this path. If
$M\le n/4$, then $n(n-1)\ge n^2/2$ and
$I\ge n^2/(16M)$. Combining this with $I\ge nd$ yields
\cref{eq:main-lower}. If $M>n/4$, then $n^2/M<4n\le2nd$,
so the output stores alone suffice. The worst-case cost is at least
the cost of this path.
\end{proof}

\section{A prefix obtained from a bilinear decomposition}
\label{sec:bilinear}
We construct a cheap prefix whose derivative field contains $\cT$.
The numerical output consists of weighted sums of scalar exponentials.
The kernels are recovered only as identities in a field.

We use the following recursive form of Strassen's identity
\cite{strassen1969}. The proof also establishes the independence
property needed for the derivative field.

\begin{lemma}[The seven-product decomposition]
\label{lem:strassen}
Let $d=2^k$ and $R=7^k=d^\sigma$, where $\sigma=\log_2 7$.
There are nonzero integer linear forms $\ell_r,m_r$ on $d\times d$
matrices and integers $c_{\alpha\beta r}$ such that
\begin{equation}
 (UZ^{\mathsf T})_{\alpha\beta}
 =\sum_{r=1}^R c_{\alpha\beta r}\ell_r(U)m_r(Z).
 \label{eq:bilinear}
\end{equation}
All $R$ forms on either input can be evaluated in $O(R)$ scalar
instructions, with an $O(R)$-transfer implementation using a constant
number of fast words. The $R$ bilinear polynomials
$\ell_r(U)m_r(Z)$ are linearly independent. The forms
$\ell_r$ span the dual of the $d^2$-dimensional input space.
\end{lemma}
\begin{proof}
The seven-product identity \cite{strassen1969}, written in
\cref{app:strassen}, gives \cref{eq:bilinear} by recursion. Its
coefficients are integers. Evaluating all leaf forms satisfies
$L(d)\le7L(d/2)+O(d^2)$, with $L(1)=O(1)$, hence $L(d)=O(R)$.
Store every intermediate in slow memory and use a constant number of
registers for each scalar instruction to obtain the stated transfer
bound.

At $d=2$, the coefficient vectors of the seven bilinear polynomials
have a $7\times7$ minor of determinant one; the minor is displayed in
\cref{app:strassen}. Recursive leaf coefficient vectors are tensor
products of these seven vectors, up to a permutation of monomials.
They are linearly independent at every depth. Finally, if all
$\ell_r(U)$ vanish, then \cref{eq:bilinear} gives $UZ^{\mathsf T}=0$
for every $Z$. Taking $Z$ to be the identity gives $U=0$. Thus the
forms $\ell_r$ span the full dual space.
\end{proof}

Assume $d\mid n$ and write $B=n/d$. Partition the rows of $Q$ and $K$
into $B$ blocks of $d$ rows, denoted $Q^{(a)}$ and $K^{(b)}$.
For $a,b\in[B]$ and $r\in[R]$, define
\begin{equation}
 u_{ar}=\ell_r(Q^{(a)}),\qquad
 z_{br}=m_r(K^{(b)}),\qquad
 E_{abr}=e^{u_{ar}z_{br}},\qquad
 t_b=V_{(b-1)d+1,1}.
 \label{eq:forms}
\end{equation}
The auxiliary vector has $BR$ entries:
\begin{equation}
 g_{ar}=\sum_{b=1}^B E_{abr}t_b.
 \label{eq:aux-vector}
\end{equation}
All original input coordinates remain free. In particular, a scalar
form $u_{ar}$ can depend on many entries of its full matrix block.

\begin{theorem}[Cost and derivative field of the auxiliary vector]
\label{thm:bilinear-prefix}
Let $d\ge8$ be a power of two, $d\mid n$, $n\ge d^2$, and $M\ge d^2$.
Computing the vector $g$ of \cref{eq:aux-vector}, with its $BR$ outputs
in designated slow addresses, has I/O complexity
\begin{equation}
 \Theta\!\left(BR+\frac{B^2R}{M}\right)
 =\Theta\!\left(nd^{\sigma-1}
              +\frac{n^2d^{\sigma-2}}M\right).
 \label{eq:aux-io}
\end{equation}
The upper bound has $O(BR+B^2R)\subseteq O(n^2d)$ work. Moreover,
\begin{equation}
 F_0(Dg)=F_0(E_{abr}:a,b\in[B],r\in[R])\supseteq\cT,
 \qquad \tau_1(g)=B^2R.
 \label{eq:aux-field}
\end{equation}
For any epoch partition of this prefix, its history at completion
contains $\cT$.
\end{theorem}
\begin{proof}
Evaluate and store all $u_{ar}$ and $z_{br}$ using
\cref{lem:strassen}. This costs $O(BR)$ work and transfers.
For a fixed $r$, divide the query and key indices into tiles of size
at most
\[
 s=\min\{B,\lfloor M/8\rfloor\}.
\]
On a tile $I\times J$, load the current accumulators $a_i$, query
scalars $u_{ir}$, key scalars $z_{jr}$, and markers $t_j$, and compute
\begin{equation}
 b_i=a_i+\sum_{j\in J}e^{u_{ir}z_{jr}}t_j\quad(i\in I).
 \label{eq:tile}
\end{equation}
At most $4s+4\le M$ fast words suffice. The tile uses
$3|I|+2|J|\le5s\le M$ transfers, including the accumulator stores.
Initialising the $B$ accumulators and visiting all tiles costs at most
\[
 B+5B\lceil B/s\rceil=O(B+B^2/M)
\]
transfers for this $r$. The total over $r$ is the upper bound in
\cref{eq:aux-io}. There are $B^2R$ scalar exponential terms and
$O(BR)$ preprocessing instructions. Since $R=d^\sigma\le d^3$,
both are within a uniform $O(n^2d)$ budget.

Fix $r$ and choose a key coordinate $\kappa_r\in[d]\times[d]$ whose
coefficient $\beta_r$ in $m_r$ is nonzero. The markers are independent
of $Q,K$. Therefore
\begin{equation}
 \frac{\partial g_{ar}}{\partial K^{(b)}_{\kappa_r}}
   =u_{ar}\beta_r E_{abr}t_b,
 \qquad
 E_{abr}=\frac{\partial g_{ar}/\partial K^{(b)}_{\kappa_r}}
                  {u_{ar}\beta_rt_b}
 \quad\text{in the function field}.
 \label{eq:extract-leaf}
\end{equation}
The denominator is a nonzero rational function, although it can
vanish on particular inputs. The prefix performs no division by it.
Every derivative of $g$ is rational in the $E_{abr}$ and the inputs,
which proves the field equality in \cref{eq:aux-field}.

For a row $i=(a-1)d+\alpha$ and $j=(b-1)d+\beta$,
\cref{eq:bilinear} gives
\begin{equation}
 W_{ij}=\prod_{r=1}^R E_{abr}^{c_{\alpha\beta r}}.
 \label{eq:kernel-product}
\end{equation}
All powers are integers, so \cref{eq:kernel-product} belongs to the
ordinary field $F_0(E_{abr}:a,b\in[B],r\in[R])$. Taking $W_{ij}/W_{in}$ gives $\cT\subseteq
F_0(E_{abr}:a,b\in[B],r\in[R])$.

For a fixed pair of blocks $(a,b)$, the exponent polynomials
$u_{ar}z_{br}$ are linearly independent by \cref{lem:strassen}.
Different block pairs use disjoint sets of query--key monomials.
Thus all $B^2R$ exponent polynomials are linearly independent and
homogeneous of degree two. \Cref{lem:independence} proves algebraic
independence of the corresponding exponentials. Hence
$\tau_1(g)=B^2R$.

The lower bound now follows from \cref{thm:history} and the $BR$
compulsory output stores. If $M\le B\sqrt R/4$, then
$I\ge B^2R/(4M)-M\ge3B^2R/(16M)$. If $M>B\sqrt R/4$, then
$B^2R/M<4B\sqrt R\le4BR$, which is controlled by the output stores.
Combining the bounds proves \cref{eq:aux-io}. Finally, the induction
in \cref{thm:history} puts $F_0(Dg)$ in the history at the end of any
partition of this executed prefix.
\end{proof}

\begin{corollary}
\label{cor:subintermediate}
Along $d=2^k\ge8$, $n=d^3$, and $M=d^2$, the prefix of
\cref{thm:bilinear-prefix} has $O(n^2d)$ work and
$\Theta(d^{\sigma+2})=o(nd+n^2d/M)$ transfers, while its history
contains $\cT$.
\end{corollary}
\begin{proof}
Here $B=d^2$, $R=d^\sigma$, and both terms of \cref{eq:aux-io}
are $d^{\sigma+2}$. In contrast, $nd+n^2d/M=d^4+d^5$.
Since $\sigma<3$, their ratio tends to zero.
\end{proof}

\section{What field containment cannot certify}
\label{sec:certificates}
The prefix of \cref{thm:bilinear-prefix} can be followed by the blocking
attention algorithm while retaining the original inputs in slow
memory. This gives a correct $O(n^2d)$-work program. The continuation
may ignore the auxiliary vector, so the argument concerns a completion
condition stated solely in terms of the history field.

\begin{theorem}[Field-containment certificates]
\label{thm:certificate}
There are no absolute constants $c,C>0$ and charges $\Pi_t$, defined
for every epoch history of every correct $O(n^2d)$-work attention
program, with the following properties: $\Pi_0=0$; a prefix history
containing $\cT$ has charge at least $c\,n(n-1)$; and every epoch
of at most $M$ transfers satisfies
\begin{equation}
 \Pi_t-\Pi_{t-1}\le C\left(M+\frac{M^2}{d}\right).
 \label{eq:putative-capacity}
\end{equation}
Charges may depend on the entire execution history and need not be
monotone.
\end{theorem}
\begin{proof}
Take $n=d^3$, $M=d^2$, and $d=2^k\ge8$. Run the prefix in
\cref{thm:bilinear-prefix}, then compute attention. Pad the end of
the prefix with fewer than $M$ transfers of constants so that it is
an epoch boundary. On other parameter triples use the blocking
algorithm; the resulting family has a uniform $O(n^2d)$ work bound.

By \cref{cor:subintermediate}, the prefix uses $O(d^\sigma)$ epochs
and its history contains $\cT$. Summing \cref{eq:putative-capacity}
from zero gives
\[
 \Pi_t\le O(d^\sigma)(d^2+d^3)
       =O(d^{\sigma+3})=o(d^6).
\]
The completion condition instead requires
$\Pi_t\ge c\,d^3(d^3-1)=\Omega(d^6)$, a contradiction.
The argument uses only upper bounds on increments.
\end{proof}

The same example excludes the smaller epoch capacity
$O(M+M^2/d^2)$. More generally, field containment alone cannot
certify a transfer cost above \cref{eq:prefix-intro} on this
parameter sequence. The additive $nd^{\sigma-1}$ term cannot be
dropped unless it is absorbed, for example when $n\ge Md$.

This is a limit on the stated completion condition, not on all
lower-bound methods. In particular, it does not address a charge
that requires numerical production of the output or exploits a
specific continuation. Nor does the prefix give an improved attention
algorithm. The uniform gap in \cref{eq:bracket} remains.

\appendix
\crefalias{section}{appendix}
\section{The bilinear identity and its independent factors}
\label{app:strassen}
For $2\times2$ matrices $U=(u_{ij})$ and $B=(b_{ij})$, use the seven
products
\begin{align*}
 p_1&=(u_{11}+u_{22})(b_{11}+b_{22}),&
 p_2&=(u_{21}+u_{22})b_{11},\\
 p_3&=u_{11}(b_{12}-b_{22}),&
 p_4&=u_{22}(b_{21}-b_{11}),\\
 p_5&=(u_{11}+u_{12})b_{22},&
 p_6&=(u_{21}-u_{11})(b_{11}+b_{12}),\\
 p_7&=(u_{12}-u_{22})(b_{21}+b_{22}).
\end{align*}
Strassen's identity \cite{strassen1969} is
\begin{equation}
 UB=\begin{pmatrix}
 p_1+p_4-p_5+p_7 & p_3+p_5\\
 p_2+p_4 & p_1-p_2+p_3+p_6
 \end{pmatrix}.
 \label{eq:strassen-base}
\end{equation}
Expanding its four entries verifies the identity over the integers.
Substitute $B=Z^{\mathsf T}$ to obtain the form used in the main text.
Applying the same identity recursively to matrix blocks gives $7^k$
leaf products for $d=2^k$. Each leaf is an integer linear form in $U$
times an integer linear form in $Z$; each output is an integer linear
combination of these products.

Order the four entries of $U$ and $B$ row by row, and order their
sixteen pairwise monomials first by the $U$ entry and then by the $B$
entry. The coefficient matrix of $(p_1,\ldots,p_7)$ has, at zero-based
column indices $0,1,3,6,8,9,12$, the minor
\begin{equation}
 \begin{pmatrix}
 1&0&1&0&0&0&1\\
 0&0&0&0&1&0&1\\
 0&1&-1&0&0&0&0\\
 0&0&0&0&0&0&-1\\
 0&0&1&0&0&0&0\\
 -1&-1&0&0&1&1&0\\
 0&0&0&1&0&0&0
 \end{pmatrix},
 \qquad \det=1.
 \label{eq:minor-seven}
\end{equation}
This proves independence at the base. At the next level, a leaf index
is a pair of leaf indices; its coefficient on a query--key monomial
is the product of the corresponding coefficients at the two levels.
After reordering the monomial coordinates, the coefficient vectors
are therefore the tensor products of the base vectors. Tensor
products of independent vectors are independent: extend the base
vectors to a basis and apply the dual basis in each tensor factor.
Induction proves independence of the $7^k$ bilinear polynomials.
Transposing the key matrix merely permutes the monomials.

The same recursion evaluates all input linear forms with
$L(d)\le7L(d/2)+c d^2$ instructions. No product of the two inputs is
needed for this preprocessing. Copies of subblocks and sums or
differences require $O(d^2)$ instructions at a node. The recurrence
sums to $O(d^\sigma)$; storing each intermediate in slow memory gives
at most a constant number of transfers per instruction. This
implementation deliberately makes no claim of optimal I/O for the
linear transforms.

\section{Numerical recovery of pairs}
\label{app:recovery}
This appendix gives a real-word counterpart of the matrix compression
step in \cite[Lemma~4.6]{sahaye2024}, with a hypothesis absent from
\cref{thm:attention}. A vector of summaries numerically recovers a
target vector when there is a continuously differentiable decoder
from those numerical values alone, on a neighbourhood of their
image. The decoder may use input-independent constants. It is not
given arbitrary elements of the rational input field.

\begin{lemma}[Rank of a selected dot-product family]
\label{lem:pair-rank}
Let $E$ be a bipartite set of $B$ selected pairs between independent
query and key vectors in $\R^d$. If $\rho$ is the generic Jacobian
rank of $(q_i\cdot k_j)_{(i,j)\in E}$, then
\begin{equation}
 B\le\frac{\rho^2}{d^2}+2\rho.
 \label{eq:pair-rank}
\end{equation}
\end{lemma}
\begin{proof}
Let
\[
 r=\sum_i\min(\deg(i),d),\qquad
 c=\sum_j\min(\deg(j),d).
\]
Choose the keys in generic position so that every set of at most $d$
of them is independent; distinct Vandermonde vectors provide a
witness. The derivatives in the query coordinates form blocks on
disjoint output coordinates, with total rank $r$. Thus $\rho\ge r$.
Interchanging the roles gives $\rho\ge c$. These generic inequalities
hold simultaneously, since the witnessing nonzero minors are
polynomials.

There are at most $r/d$ query vertices and $c/d$ key vertices with
degree at least $d$. At most $rc/d^2$ edges join two such vertices.
Edges incident to a lower-degree query number at most $r$, and the
remaining edges incident to a lower-degree key number at most $c$.
Hence $B\le rc/d^2+r+c\le\rho^2/d^2+2\rho$.
\end{proof}

\begin{theorem}[Numerical compression]
\label{thm:compression}
Suppose $m$ continuously differentiable real summaries numerically
recover $B$ selected entries from one of the families
\[
 q_i\cdot k_j,\qquad e^{q_i\cdot k_j},\qquad
 e^{q_i\cdot(k_j-k_n)}\ (j<n)
\]
on a nonempty open input set. Then
\begin{equation}
 B\le\frac{m^2}{d^2}+2m.
 \label{eq:compression}
\end{equation}
The summaries may depend on $V$ and on every input coordinate.
\end{theorem}
\begin{proof}
Numerical factorisation bounds the rank of the recovered vector by
$m$, by \cref{lem:factor-rank}. Apply \cref{lem:pair-rank} for scores.
Coordinate-wise exponentiation multiplies its Jacobian by an
invertible diagonal matrix, preserving rank. For the centred family,
change variables from $K$ to $(K_1-K_n,\ldots,K_{n-1}-K_n,K_n)$.
This is an invertible linear change, and the centred key vectors are
independent coordinates. The selected minors are nonzero on a dense
open subset, so the argument applies inside the stated open input
set.
\end{proof}

\begin{corollary}[The full bound under pair coverage]
\label{cor:pair-cover}
In the attention parameter range, suppose the epochs of an open
execution path numerically recover, between them, every entry of one
full family in \cref{thm:compression}. The path has
$\Omega(nd+n^2d^2/M)$ transfers, with an absolute constant.
\end{corollary}
\begin{proof}
Every value computed during an epoch is an analytic function of its
at most $2M$ incoming values. By \cref{thm:compression}, one epoch
recovers at most $4M^2/d^2+4M\le8M^2/d^2$ selected pairs.
There are at least $n(n-1)$ distinct targets in a full family.
Together with \cref{eq:epochs}, this gives
\[
 I\ge\frac{n(n-1)d^2}{8M}-M.
\]
If $M\le nd/8$, the right side is at least a constant times
$n^2d^2/M$. Otherwise this latter term is at most $8nd$.
The $nd$ output stores complete the bound in both cases.
\end{proof}

A program can compute $W_{ij}$ as
$\prod_{\ell=1}^d e^{Q_{i\ell}K_{j\ell}}$, without forming the sum
$q_i\cdot k_j$. Its completed kernels satisfy the hypothesis of
\cref{cor:pair-cover}. This example prevents interpreting pair
coverage as an explicit-score assumption. No coverage property is
proved for every exact program.

\section{The blocking upper bound and small parameter cases}
\label{app:upper}
We include the schedule only to specify the uniform comparison in
\cref{eq:bracket}; the bound is the established attention tiling bound
\cite{dao2022,sahaye2024}. For $M\ge16d$, choose
\[
 b=\min\{n,\lfloor M/(8d)\rfloor\}.
\]
Keep $b$ query rows, $b$ length-$d$ numerator accumulators, and $b$
denominators in fast memory. Scan the key--value rows. For each such
row, compute the $b$ exponential weights and update the accumulators.
All denominators are positive. Divide the numerators and write the
output block, then proceed to the next query block.

The query block and numerators use $2bd$ words. Denominators use $b$,
a key and value row use $2d$, and eight scratch words suffice for the
scalar inner loops. For $d\ge2$ and $M\ge16d$,
\[
 2bd+b+2d+8\le M.
\]
The transfers are $O(nd+\lceil n/b\rceil nd)$, hence
$O(nd+n^2d^2/M)$. Every pair is processed with $O(d)$ work.
A product of coordinate-wise exponentials can replace the dot-product
sum with the same asymptotic work and transfer count.

If $d^2\le M<16d$, then $d<16$ and $M<256$. A scalar implementation
can spill all nonresident partial sums and evaluate the attention
formula with four fast words and $O(n^2d+nd)$ transfers. For example,
one register holds the current sum, two hold the next input operands,
and one holds a temporary product or exponential; completed values
and outer accumulators may be spilled before the next step. Since
$M/d<16$, this count is within an absolute constant of
$nd+n^2d^2/M$. The work remains $O(n^2d)$. This covers the bounded
parameter cases as well as the regime in which full rows fit.

\section*{AI methodology}
\phantomsection
\addcontentsline{toc}{section}{AI methodology}
The proofs in this paper were developed with AI assistance from GPT-6
Astra Pro and Claude models and were checked line by line by the author.

\begingroup
\microtypesetup{protrusion=false,expansion=false}
\urlstyle{same}
\bibliographystyle{alphaurl}
\bibliography{references}
\endgroup
\end{document}
